\documentclass[border=10pt]{standalone}
\usepackage[utf8]{inputenc}
\usepackage[spanish]{babel}
\usepackage{tikz}
\usetikzlibrary{shapes.geometric, arrows.meta, positioning, calc}

\begin{document}

% Definición de estilos
\tikzset{
    % Cajas de proceso (rectángulos con bordes redondeados)
    proceso/.style={
        rectangle,
        rounded corners=3pt,
        draw=cyan!70!black,
        fill=cyan!10,
        text width=4cm,
        minimum height=1.5cm,
        align=center,
        font=\small
    },
    % Caja de resultado final
    resultado/.style={
        rectangle,
        rounded corners=3pt,
        draw=cyan!70!black,
        fill=cyan!10,
        text width=2.5cm,
        minimum height=1.5cm,
        align=center,
        font=\small
    },
    % Círculos numerados (pasos)
    paso/.style={
        circle,
        fill=orange,
        text=white,
        font=\bfseries\small,
        minimum size=0.7cm,
        inner sep=0pt
    },
    % Flechas
    flecha/.style={
        ->,
        >=Stealth,
        thick,
        cyan!70!black
    }
}

\begin{tikzpicture}[node distance=0.8cm and 1.5cm]

% === FILA SUPERIOR (Con reserva) ===
% Paso 1 superior
\node[paso] (p1a) {1};
\node[proceso, right=0.3cm of p1a] (caja1a) {Sumar la cantidad de personas que entraron con reserva durante la semana.};

% Paso 2 superior
\node[paso, right=1.5cm of caja1a] (p2a) {2};
\node[proceso, right=0.3cm of p2a] (caja2a) {Multiplicar la cantidad obtenida en el paso \tikz[baseline=-0.5ex]{\node[circle,fill=orange,text=white,font=\bfseries\tiny,inner sep=1pt]{1};} por 22,5.};

% === FILA INFERIOR (Sin reserva) ===
% Paso 1 inferior
\node[proceso, below=1.5cm of caja1a] (caja1b) {Sumar la cantidad de personas que entraron sin reserva durante la semana.};
\node[paso, left=0.3cm of caja1b] (p1b) {1};

% Paso 2 inferior
\node[proceso, right=1.5cm of caja1b] (caja2b) {Multiplicar la cantidad obtenida en el paso \tikz[baseline=-0.5ex]{\node[circle,fill=orange,text=white,font=\bfseries\tiny,inner sep=1pt]{1};} por 17.};
\node[paso, left=0.3cm of caja2b] (p2b) {2};

% === PASO 3 (Resultado) ===
\node[paso, right=1.5cm of caja2a, yshift=-1.5cm] (p3) {3};
\node[resultado, right=0.3cm of p3] (caja3) {Comparar los resultados};

% === FLECHAS ===
% Flechas horizontales fila superior
\draw[flecha] (caja1a.east) -- (p2a.west);
\draw[flecha] (caja2a.east) -- ++(0.5,0) |- (caja3.west);

% Flechas horizontales fila inferior
\draw[flecha] (caja1b.east) -- (p2b.west);
\draw[flecha] (caja2b.east) -- ++(0.5,0) |- (caja3.west);

\end{tikzpicture}

\end{document}
